% Pattern: head/case. The steps of the pf style of proofs, \begin{step+},
% have no counter of their own. When a file asks for numbers, as
% \pfmixednumbers does, a step is numbered by its place in its proof,
% after the number of the step that the proof proves, as 1.2 for the second
% step of the proof of step 1. Seen in the gold pilot (set X).
% want theorem-1 statement/theorem 1
% want proof-1 proof - proves=theorem-1
% want step--1 other/case 1
% want proof-2 proof -
% want step--2 other/case 1.1
% want step--3 other/case 1.2
% want step--4 other/case 2
\documentclass{article}
\newtheorem{theorem}{Theorem}
\newcommand{\pfmixednumbers}[1]{}
\pfmixednumbers{4}
\newenvironment{step+}[1]{\par\textbf{Step.} #1\par}{\par}
\begin{document}
\begin{theorem}
Every bounded set is small.
\end{theorem}
\begin{proof}
\begin{step+}{The set has a largest point.}
\begin{proof}
\begin{step+}{The set is not empty.}
It has a point by the choice of the set.
\end{step+}
\begin{step+}{The points are bounded.}
The bound holds for each point.
\end{step+}
\end{proof}
\end{step+}
\begin{step+}{The largest point bounds the set.}
This is the definition of a bound.
\end{step+}
\end{proof}
\end{document}
